\bookchapter{Capitolo 62: Sviluppare l'interfaccia — introduzione a SwiftUI}

\chapterintro{L'interfaccia è tutto ciò che l'utente vede. SwiftUI è il framework dichiarativo presentato da Apple nel 2019: ti basta “dichiarare” come deve apparire l'interfaccia e il sistema si occupa da solo di disegnarla, adattarla alla modalità scura e ai vari schermi. Nel 2026 SwiftUI è ormai lo standard assoluto dei nuovi progetti. In questo capitolo impariamo a “disegnare”.}

\booksection{L'idea di SwiftUI: dichiarativo, non imperativo}

Lo sviluppo delle interfacce all'antica è “imperativo”: dici al sistema “crea il pulsante mettilo a (20,30) imposta il colore”, una istruzione alla volta. SwiftUI è “dichiarativo”: descrivi direttamente “sull'interfaccia c'è un pulsante rosso con scritto 'Toccami'”, e come disegnarlo e quando ridisegnarlo se lo pensa il sistema.

\begin{lstlisting}[language=Swift]
import SwiftUI

struct ContentView: View {
    var body: some View {      // body: descrive "come appare lo schermo"
        VStack {               // disposizione verticale
            Text("Ciao, mondo!")
                .font(.largeTitle)      // modificatore: titolo grande
            Button("Toccami") {
                print("cliccato")
            }
            .buttonStyle(.borderedProminent)  // pulsante blu pieno
        }
        .padding()
    }
}
\end{lstlisting}

Punti chiave: una vista è uno struct conforme al protocollo View, la proprietà body descrive l'interfaccia, e i .modificatori() si concatenano a catena come una “truccatura”. Questo è tutto il fondamento di SwiftUI.

\booksection{Le viste di base: testo, pulsanti, campi di input, immagini}

\begin{center}
\begin{tabularx}{\linewidth}{|X|X|X|}
\hline
\textbf{Vista} & \textbf{A cosa serve} & \textbf{Esempio} \tabularnewline
\hline
Text & Mostrare testo & Text("Ciao") \tabularnewline
Button & Pulsante & Button("Toccami") \{ azione \} \tabularnewline
TextField & Campo di input & TextField("Segnaposto", text: \$testo) \tabularnewline
Image & Mostrare un'immagine & \texttt{Image(systemName: "star.\allowbreak fill")} \tabularnewline
\hline
\end{tabularx}
\end{center}

\begin{center}\small\color{gray}Le quattro viste fondamentali di SwiftUI\end{center}

Nota bene Image(systemName:): Apple include la libreria di icone SF Symbols (migliaia di icone gratuite) e basta un nome come “star.fill” per mostrare all'istante una stella piena. È un vantaggio nascosto dello sviluppo iOS: per costruire interfacce quasi non serve cercare materiale grafico.

\booksection{Il trio del layout: VStack, HStack e ZStack}

\begin{itemize}
  \item VStack: dispone in verticale (come la fila alla mensa);
  \item HStack: dispone in orizzontale (come sedersi fianco a fianco);
  \item ZStack: sovrappone (come carte da gioco impilate, ideale per badge e coperture).
\end{itemize}

\begin{lstlisting}[language=Swift]
VStack(alignment: .leading, spacing: 12) {
    HStack {
        Text("Conto")
        Spacer()                 // spazio elastico: spinge la parte destra al bordo
        Text("-25.5 yuan").foregroundColor(.red)
    }
    Divider()                    // divisore
    Text("Nota: pranzo")
}
.padding(16)                     // margine su tutti i lati
.background(Color(.systemGray6)) // sfondo grigio chiaro
.cornerRadius(12)                // angoli arrotondati

// L'ordine dei modificatori conta: prima padding poi background, lo sfondo non include il margine
\end{lstlisting}

\begin{tcbwarn}{Attenzione: l'ordine dei modificatori conta}
padding(16).background(grigio) e background(grigio).padding(16) danno risultati diversi: nel primo caso “prima lo spazio, poi il colore” (lo sfondo non include i margini), nel secondo “prima il colore, poi lo spazio” (lo sfondo include i margini). Da qui nasce il “bordo bianco misterioso” più comune tra i principianti. Se hai dubbi, scambia l'ordine: l'anteprima ti dà subito la risposta.

\end{tcbwarn}

\booksection{Stato e binding: far “vivere” l'interfaccia}

Un'interfaccia statica è noiosa. La regola di SwiftUI: l'interfaccia rispecchia sempre fedelmente lo stato; quando lo stato cambia, l'interfaccia si ridisegna da sola.

\begin{lstlisting}[language=Swift]
struct CounterView: View {
    @State private var count = 0   // @State: se questo valore cambia, ridisegna

    var body: some View {
        VStack {
            Text("Hai toccato \(count) volte")
            Button("Toccami +1") {
                count += 1          // cambia lo stato; l'interfaccia si aggiorna da sola
            }
        }
    }
}
\end{lstlisting}

I campi di input usano il “binding bidirezionale”: il simbolo \$ passa insieme “il valore e il diritto di modificarlo”:

\begin{lstlisting}[language=Swift]
struct InputView: View {
    @State private var text = ""

    var body: some View {
        TextField("Scrivi qualcosa", text: $text)  // $ binding bidirezionale
            .textFieldStyle(.roundedBorder)
        Text("Hai digitato: \(text)")   // si aggiorna in tempo reale
    }
}
\end{lstlisting}

Azione dell'utente

Lo stato cambia (count += 1)

SwiftUI rileva il cambiamento

L'interfaccia si ridisegna da sola (body ricalcolato)

\begin{center}\small\color{gray}Figura 3: il flusso di dati a senso unico “stato interfaccia” di SwiftUI\end{center}

\booksection{Anteprime, tema e animazioni: i “potenziamenti” di SwiftUI}

\begin{itemize}
  \item La macro \#Preview: aggiungi in fondo al file \#Preview \{ ContentView() \} e il canvas a destra esegue il rendering in tempo reale: modifichi il codice e vedi subito l'effetto;
  \item Modalità scura gratis: usa i colori di sistema (Color(.systemGray6), .primary) invece dei colori fissi e l'adattamento alla modalità scura avviene da solo, senza scrivere una riga in più;
  \item Animazioni in una riga: withAnimation \{ count += 1 \} oppure aggiungi .animation(.default, value: count) alla vista: transizioni scorrevoli al punto giusto.
\end{itemize}

\begin{lstlisting}[language=Swift]
struct AnimatedHeartView: View {
    @State private var liked = false

    var body: some View {
        Button {
            withAnimation(.spring(duration: 0.4)) {
                liked.toggle()
            }
        } label: {
            Image(systemName: liked ? "heart.fill" : "heart")
                .font(.system(size: 48))
                .foregroundColor(liked ? .red : .gray)
                .scaleEffect(liked ? 1.3 : 1.0)
        }
    }
}
\end{lstlisting}

\begin{tcbinfo}{Consiglio: prima le funzioni, poi le animazioni}
Il difetto tipico dei principianti è “mettersi a dibattere le animazioni prima che l'interfaccia sia giusta”. Fai prima funzionare tutto e aggiungi gli effetti alla fine, in un colpo solo: l'animazione è la ciliegina sulla torta, la funzione è il pane quotidiano. È anche il nostro principio di sempre: “la semplicità prima di tutto”.

\end{tcbinfo}

\begin{tcbpractice}{Mettiti alla prova: crea un "pulsante dell'umore"}
\booksection{Scheda rapida dei modificatori più usati: una pagina a portata di mano e l'interfaccia non è più un problema}

\begin{center}
\begin{tabularx}{\linewidth}{|X|X|}
\hline
\textbf{Cosa vuoi ottenere} & \textbf{Modificatore} \tabularnewline
\hline
Spaziatura & \texttt{.\allowbreak padding(16)} \tabularnewline
Angoli arrotondati & \texttt{.\allowbreak cornerRadius(12)} \tabularnewline
Colore di sfondo & \texttt{.\allowbreak background(Color(.\allowbreak systemGray6))} \tabularnewline
Carattere e corpo & \texttt{.\allowbreak font(.\allowbreak title.\allowbreak bold())}/\texttt{.\allowbreak font(.\allowbreak system(size: 20))} \tabularnewline
Colore in primo piano (testo/icone) & \texttt{.\allowbreak foregroundColor(.\allowbreak red)} \tabularnewline
Riempi la larghezza & \texttt{.\allowbreak frame(maxWidth: .\allowbreak infinity)} \tabularnewline
Dimensione fissa & \texttt{.\allowbreak frame(width: 100, height: 44)} \tabularnewline
Allineamento & \texttt{VStack(alignment: .\allowbreak leading)} \tabularnewline
Trasparenza & \texttt{.\allowbreak opacity(0.\allowbreak 6)} \tabularnewline
Ombra & \texttt{.\allowbreak shadow(radius: 5)} \tabularnewline
Gesto di tocco & \texttt{.\allowbreak onTapGesture \{ .\allowbreak .\allowbreak .\allowbreak  \}} \tabularnewline
Tipo di tastiera & \texttt{.\allowbreak keyboardType(.\allowbreak decimalPad)} \tabularnewline
\hline
\end{tabularx}
\end{center}

\begin{center}\small\color{gray}Scheda rapida dei modificatori più usati di SwiftUI\end{center}

Se non li ricordi non importa: questa pagina serve proprio da consultazione. In realtà i più usati sono i primi quattro: padding, angoli arrotondati, sfondo e carattere.

Trasforma il progetto modello in un selettore d'umore: un pulsante che a ogni tocco cicla tra felice triste arrabbiato, con testo e colore di sfondo che cambiano di conseguenza. Suggerimento: salva in un array le tre coppie (testo+colore), tieni l'indice in un @State e al tocco usa (index + 1) \% 3. A lavoro finito mandalo agli amici e raccogli la prima ondata di “Wow, sai creare app!”.

\end{tcbpractice}

\begin{tcbsummary}{Riepilogo del capitolo}
\begin{itemize}
  \item SwiftUI è dichiarativo: descrivi "come deve apparire" e il sistema si occupa di disegnare e ridisegnare
  \item Vista = uno struct conforme al protocollo View; body descrive l'interfaccia; i modificatori si concatenano come un "trucco"
  \item Le quattro viste di base: Text / Button / TextField / Image; la libreria di icone SF Symbols è gratuita e in abbondanza
  \item Il trio del layout: VStack in verticale, HStack in orizzontale, ZStack in sovrapposizione; Spacer è lo spazio elastico
  \item @State dà vita all'interfaccia; il \$ del binding bidirezionale alimenta i campi di input
  \item \#Preview fa vedere subito il risultato; i colori di sistema si adattano da soli alla modalità scura; withAnimation aggiunge effetti in una riga
\end{itemize}

\end{tcbsummary}
